\documentclass[11pt]{article}
\usepackage[T1]{fontenc}
\usepackage{lmodern}
\usepackage[margin=1in]{geometry}
\usepackage{amsmath,amssymb,amsthm,mathtools}
\usepackage{microtype}
\usepackage{xcolor}
\usepackage{enumitem}
\PassOptionsToPackage{hyphens}{url}
\usepackage{hyperref}
\usepackage{bookmark}
\hypersetup{colorlinks=true,linkcolor=blue!45!black,citecolor=blue!45!black,
  urlcolor=blue!45!black,
  pdftitle={A Counterexample to the Group-Ring Determinant Conjecture}}
\IfFileExists{glyphtounicode-cmex.tex}{\input{glyphtounicode-cmex.tex}}{}
\pdfgentounicode=1
\newtheorem{theorem}{Theorem}[section]
\newtheorem{proposition}[theorem]{Proposition}
\newtheorem{lemma}[theorem]{Lemma}
\theoremstyle{remark}
\newtheorem{remark}[theorem]{Remark}
\newtheorem{example}[theorem]{Example}
\numberwithin{equation}{section}
\newcommand{\F}{\mathbb F}
\newcommand{\Z}{\mathbb Z}
\newcommand{\Q}{\mathbb Q}
\newcommand{\R}{\mathbb R}
\newcommand{\C}{\mathbb C}
\newcommand{\Mat}{\operatorname{M}}
\newcommand{\diag}{\operatorname{diag}}
\newcommand{\N}{\mathcal N}
\newcommand{\unverified}[1]{\par\smallskip\noindent\fbox{\parbox{0.95\linewidth}{\textbf{Unverified step.} #1}}\par\smallskip}
\setlength{\emergencystretch}{2em}
\title{A Counterexample to the Group-Ring Determinant Conjecture}
\author{Rewritten for readability from the OpenAI manuscript of the same title}
\date{September 23, 2026}

\begin{document}
\maketitle

\begin{abstract}
We disprove the unrestricted group-ring Determinant Conjecture. We construct a finitely generated group $G$ and a square matrix over $\Z[G]$ that is invertible over $\Q[G]$ and has Fuglede--Kadison determinant strictly between zero and one. The logarithmic integral defining the determinant is finite. The construction rests on one external input (Theorem~\ref{thm:input}); everything after that input is proved here.
\end{abstract}

\section{The result}
\label{sec:result}

\subsection{Conventions}

Throughout, $G$ is a discrete group. For a commutative ring $R$, $R[G]$ is the group ring of \emph{finite} sums $c=\sum_g c_g g$ with $c_g\in R$.

\begin{itemize}[leftmargin=2em]
\item \emph{Identity coefficient.} $[1]c:=c_1$, the coefficient of the identity element.
\item \emph{Involution.} For $R\subseteq\C$, $c^*:=\sum_g\overline{c_g}\,g^{-1}$. For a matrix $A$ over $R[G]$, $(A^*)_{ij}:=(A_{ji})^*$.
\item \emph{Projection.} A matrix $P\in\Mat_m(\Q[G])$ is an \emph{orthogonal projection} if $P=P^*=P^2$.
\item \emph{Regular representation.} $G$ acts on $\ell^2(G)$ by the unitaries $(\lambda(g)\xi)(h)=\xi(g^{-1}h)$, so $\lambda(g)\delta_h=\delta_{gh}$. Extend $\lambda$ linearly to $\C[G]$. For $A\in\Mat_n(\C[G])$, the operator $T_A$ on $\ell^2(G)^n$ is $(T_A\xi)_i=\sum_j\lambda(A_{ij})\xi_j$. Since $\lambda$ is a unitary representation, $A\mapsto T_A$ is a unital $*$-homomorphism: $T_{AB}=T_AT_B$ and $T_{A^*}=T_A^*$.
\item \emph{Unit vectors.} $\varepsilon_i\in\ell^2(G)^n$ has $\delta_1$ in component $i$ and $0$ elsewhere. (The source calls these $e_i$; they are renamed because $e$ is used below for a reduction of $P$.) Since $\lambda(g)\delta_1=\delta_g$,
\begin{equation}\label{eq:trace-coefficient}
 \langle T_A\varepsilon_i,\varepsilon_i\rangle=[1]A_{ii}.
\end{equation}
\item \emph{Trace.} For a bounded operator $T$ on $\ell^2(G)^n$, $\tau_n(T):=\sum_{i=1}^n\langle T\varepsilon_i,\varepsilon_i\rangle$. This trace is \emph{unnormalized}: $\tau_n(I_n)=n$. For a group-ring matrix we write $\tau_n(A):=\tau_n(T_A)=\sum_i[1]A_{ii}$, by Equation~\eqref{eq:trace-coefficient}.
\end{itemize}

\subsection{The conjecture}

For $A\in\Mat_n(\Z[G])$, the probability spectral measure of the positive operator $T_A^*T_A$ is
\[
 \mu_A(E)=\frac1n\sum_{i=1}^n
 \langle \mathbf1_E(T_A^*T_A)\varepsilon_i,\varepsilon_i\rangle
 \quad(E\subseteq[0,\infty)\text{ Borel}).
\]
The square-matrix instance of the Determinant Conjecture is the assertion
\begin{equation}\label{eq:conjecture}
 \int_{(0,\infty)}\log t\,d\mu_A(t)\ge0\qquad\text{for all }G,\ n,\ A\in\Mat_n(\Z[G]).
\end{equation}
The zero atom is omitted, and in general the integral may be $-\infty$. The usual conjecture also permits rectangular matrices \cite[Conjecture 13.2]{Luck2002}; a square counterexample already contradicts it.

Write $\N(G)$ for the group von Neumann algebra. If $T_A$ is invertible, then $\|T_A^{-1}\|^{-2}I\le T_A^*T_A\le\|T_A\|^2I$, so $\mu_A$ is supported in a compact subinterval of $(0,\infty)$, on which $\log$ is continuous. The spectral theorem then gives $\langle\log(T_A^*T_A)\varepsilon_i,\varepsilon_i\rangle=\int\log t\,d\langle\mathbf 1_{(\cdot)}(T_A^*T_A)\varepsilon_i,\varepsilon_i\rangle$. Summing over $i$, the determinant convention for invertible $T_A$ is
\begin{equation}\label{eq:normalization}
 \log\det_{\N(G)}(T_A)
 :=\frac12\tau_n\bigl(\log(T_A^*T_A)\bigr)
 =\frac n2\int_{(0,\infty)}\log t\,d\mu_A(t).
\end{equation}
For the trivial group this is $\log|\det A|$. This determinant was introduced by Fuglede and Kadison~\cite{FK1952}.

\begin{remark}[Matching the row-vector convention]\label{rem:convention}
The inversion unitary $\delta_g\mapsto\delta_{g^{-1}}$ on each copy of $\ell^2(G)$ turns left multiplication by $g$ into right multiplication by $g^{-1}$. It therefore identifies $T_A$ with right multiplication by $A^*$ on row vectors. It fixes $\delta_1$, hence every $\varepsilon_i$, and so preserves the spectral measures above. Our convention and the row-vector convention thus state the same conjecture.
\end{remark}

\subsection{Main theorem and its one external input}

\begin{theorem}\label{thm:main}
There are a finitely generated discrete group $G_D$, an integer
$n\ge1$, and $A\in\Mat_n(\Z[G_D])$ invertible over $\Q[G_D]$ such
that
\[
 0<\det_{\N(G_D)}(T_A)<1.
\]
The integral in Equation~\eqref{eq:conjecture} is finite and strictly
negative. Thus the unrestricted group-ring Determinant Conjecture
has a negative resolution.
\end{theorem}

A unital ring is \emph{directly finite} if $ab=1$ implies $ba=1$. The proof uses the following scalar failure of direct finiteness.

\begin{theorem}[{\cite[Theorem 1.1]{CharacteristicTwo}}]
\label{thm:input}
There exist a finite field $K_2$ of characteristic two, a finitely
generated group $H_2$, and finite sums $a_2,b_2\in K_2[H_2]$ with
$a_2b_2=1$ and $b_2a_2\ne1$.
\end{theorem}

\unverified{Theorem~\ref{thm:input} is an external existence statement that is not proved in the source or in this document, and it cannot be checked from either. It asserts a failure of direct finiteness for a group algebra over a finite field, that is, a counterexample to Kaplansky's direct finiteness conjecture. Combined with the sofic determinant bound of Elek and Szab\'o \cite[Theorem 5]{ElekSzabo2005}, it also yields the existence of a non-sofic group. Both are, as far as we know, open problems, so the claim is extraordinary. Everything downstream (Lemma~\ref{lem:binary} through Theorem~\ref{thm:main}) is conditional on it.}

\begin{remark}[Status of the input]\label{rem:input}
Theorem~\ref{thm:input} is quoted, not proved or checked, in this document. It is a strong statement, and Theorem~\ref{thm:main} is exactly as reliable as it is. According to the source, Section~6 of the companion specifies the field, the group presentation and the scalar sums by a terminating finite prescription. Only the existence statement is used here, not an executed numerical instance or a group-word decision procedure. The field $K_2$ is finite but is not asserted to be $\F_2$. There is no torsion-free hypothesis, and the companion construction uses torsion. Characteristic two is essential below: a result only in odd characteristic would not supply the binary premise.
\end{remark}

\section{Proof at a glance}
\label{sec:glance}

\paragraph{Key idea.}
Call a triple $(P,D,(u,v))$ over a group $G$ a \emph{gap datum of sizes $(m,r)$} if
\begin{itemize}[leftmargin=2em]
\item $P=P^*=P^2\in\Mat_m(\Q[G])$;
\item $D$ is an odd positive integer with $DP\in\Mat_m(\Z[G])$;
\item the reduction $\bar P$ of $P$ modulo two factors as $\bar P=uv$, where $u$ is $m\times r$ and $v$ is $r\times m$ over $\F_2[G]$.
\end{itemize}
The proof has two independent halves.
\begin{enumerate}[leftmargin=2em]
\item \emph{A general machine} (Sections~\ref{sec:lifting} and~\ref{sec:analysis}). Every gap datum yields an integral matrix $A$ of size $r+m$, invertible over $\Q[G]$, with
\[
 \log\det_{\N(G)}(T_A)=(r-\tau_m(P))\log q,
\]
where $q>1$ is a power of two with $q\equiv1\pmod D$. This is negative as soon as $\tau_m(P)>r$.
\item \emph{A gap datum with $\tau_m(P)>r$} (Section~\ref{sec:projection}). From the input $a_2b_2=1\ne b_2a_2$ we build a gap datum of sizes $(d+1,d)$ over $G_D=(\bigoplus_{h\in H_2}C_3)\rtimes H_2$ with $\tau_{d+1}(P)=d+\theta$, $\theta>0$.
\end{enumerate}

\paragraph{Steps and dependencies.}
\begin{itemize}[leftmargin=2em]
\item \textbf{Lemma~\ref{lem:binary}} (uses Theorem~\ref{thm:input}). Replace the field $K_2$ by $\F_2$ at the cost of $d\times d$ matrices: $X,Y\in\Mat_d(\F_2[H_2])$ with $XY=I_d$ and a nonzero idempotent defect $C=I_d-YX$ satisfying $XC=0$, $CY=0$.
\item \textbf{Proposition~\ref{prop:separation}} (uses Lemma~\ref{lem:binary}). Some coefficient of some entry of $C$ equals $1$. Averaging projections over the $C_3$ coordinates produce a rational projection $p$ of trace $\theta>0$ that isolates this coefficient modulo two. Then $P=\diag(I_d,p)$ has trace $d+\theta$, while $\bar P$ factors through $d$ columns.
\item \textbf{Lemma~\ref{lem:lift}} (self-contained). The factorization of $\bar P$ modulo $2$ lifts to a factorization of $P$ modulo $q=2^k$, by a finite geometric series.
\item \textbf{Lemma~\ref{lem:integral}} (uses Lemma~\ref{lem:lift}). With integral lifts $U,V$ of the mod-$q$ factors, an explicit block matrix $A$ is integral, and $A=L_1BR_1$ for block shears $L_1,R_1$ and a block-diagonal $B$ built from $q$ and $P$. Hence $A$ is invertible over $\Q[G]$.
\item \textbf{Lemma~\ref{lem:logderivative}} (self-contained analysis). $\frac{d}{dx}\tau_n(\log(M^*M))=2\operatorname{Re}\tau_n(M'M^{-1})$.
\item \textbf{Lemma~\ref{lem:shears}} (uses Lemma~\ref{lem:logderivative}). Along the path $M(x)=L_xBR_x$ the quantity $\tau(\log(M^*M))$ is constant.
\item \textbf{Proposition~\ref{prop:criterion}} (uses Lemmas~\ref{lem:integral} and~\ref{lem:shears}). Hence $\tau(\log(T_A^*T_A))=\tau(\log(B^*B))=2(r-\tau_m(P))\log q$.
\item \textbf{Theorem~\ref{thm:main}} (uses Propositions~\ref{prop:separation} and~\ref{prop:criterion}). Insert $r=d$, $\tau_{d+1}(P)=d+\theta$ to get $\det_{\N(G_D)}(T_A)=q^{-\theta}$.
\end{itemize}
The projection construction, the modular lifting and the trace calculation have distinct roles; none replaces the scalar existence input.

\paragraph{Renamed symbols.}
Relative to the source, a few symbols that had two meanings are renamed.
\begin{itemize}[leftmargin=2em]
\item The unit vectors are $\varepsilon_i$, and the regular representation is $\lambda(g)$.
\item The coefficient ring $(\Z/q\Z)[G]$ is $\mathcal R_q$.
\item In Lemma~\ref{lem:logderivative}, the bounds are $c_0,c_1$, the contraction is $\Lambda$, and its norm bound is $\eta$.
\item In Lemma~\ref{lem:shears}, the shear blocks are $\Gamma,\Delta$.
\item Generic elements in the trace identities are $\alpha,\beta$.
\end{itemize}

\section{From the scalar defect to a projection}
\label{sec:projection}

This section produces the gap datum. All group-ring sums are finite, and matrix products are written in the order in which they are multiplied.

\begin{lemma}\label{lem:binary}
Let $K_2$ be a finite field of characteristic two and $H_2$ a group.
For $d=[K_2:\F_2]$, there is an injective unital ring homomorphism
\[
 \rho:K_2[H_2]\longrightarrow\Mat_d(\F_2[H_2]).
\]
For $a_2,b_2\in K_2[H_2]$ with $a_2b_2=1\ne b_2a_2$, their images $X=\rho(a_2)$, $Y=\rho(b_2)$ satisfy
$XY=I_d$ and $C:=I_d-YX\ne0$, with
\begin{equation}\label{eq:defect}
 C^2=C,\qquad XC=0,\qquad CY=0.
\end{equation}
\end{lemma}

This lemma is used once, to feed Proposition~\ref{prop:separation}.

\begin{proof}
\emph{Step 1 (the field).} Choose an $\F_2$-basis of $K_2$. For $\kappa\in K_2$ let $m(\kappa)\in\Mat_d(\F_2)$ be the matrix of the $\F_2$-linear map $t\mapsto\kappa t$. Then $m$ is additive, $m(1)=I_d$ and $m(\kappa\kappa')=m(\kappa)m(\kappa')$. It is injective: if $m(\kappa)=0$ then $\kappa=\kappa\cdot1=0$.

\emph{Step 2 (the group ring).} A matrix over $\F_2[H_2]$ is the same thing as a finite formal sum $\sum_hM_hh$ with $M_h\in\Mat_d(\F_2)$, and scalar matrices commute with group elements. Define $\rho(\sum_h\kappa_hh)=\sum_hm(\kappa_h)h$. It is additive and unital, and
\[
 \rho\Bigl(\sum_{g,h}\kappa_g\kappa'_h\,gh\Bigr)
 =\sum_{g,h}m(\kappa_g)m(\kappa'_h)\,gh
 =\Bigl(\sum_gm(\kappa_g)g\Bigr)\Bigl(\sum_hm(\kappa'_h)h\Bigr),
\]
so $\rho$ is multiplicative. If $\rho(\sum_h\kappa_hh)=0$, then $m(\kappa_h)=0$ for each $h$, so each $\kappa_h=0$ by Step~1. Thus $\rho$ is injective.

\emph{Step 3 (the identities).} $XY=\rho(a_2b_2)=I_d$. By injectivity, $YX=\rho(b_2a_2)\ne\rho(1)=I_d$, so $C\ne0$. Using $XY=I_d$:
\begin{align*}
 (YX)^2&=Y(XY)X=YX,\ \text{hence } C^2=I_d-2YX+(YX)^2=C;\\
 XC&=X-(XY)X=0;\\
 CY&=Y-Y(XY)=0.\qedhere
\end{align*}
\end{proof}

\paragraph{The group $G_D$.}
Let $H_2$ be a group. In $\bigoplus_{h\in H_2}C_3$, let $z_h$ generate the copy of $C_3$ labelled by $h$; distinct $z_h$ commute. $H_2$ acts by left translation of the labels, and
\[
 G_D=\Bigl(\bigoplus_{h\in H_2}C_3\Bigr)\rtimes H_2,
 \qquad g\,z_h\,g^{-1}=z_{gh}\quad(g,h\in H_2).
\]
Since $H_2\le G_D$, we have $\F_2[H_2]\subseteq\F_2[G_D]$.

\begin{proposition}\label{prop:separation}
Let $X,Y\in\Mat_d(\F_2[H_2])$, $d\ge1$, satisfy
$XY=I_d\ne YX$. In the group
$G_D=(\bigoplus_{h\in H_2}C_3)\rtimes H_2$ above, there is an orthogonal
projection $P\in\Mat_{d+1}(\Z[1/3][G_D])$ and a finite set
$S\subseteq H_2\setminus\{1\}$ such that
\[
 \tau_{d+1}(P)=d+\theta,\qquad
 \theta=\frac13\left(\frac23\right)^{|S|}>0,
 \qquad \bar P=uv,
\]
where $u$ has size $(d+1)\times d$ and $v$ has size $d\times(d+1)$
over $\F_2[G_D]$. The integer $D=3^{1+|S|}$ clears all coefficients
of $P$.
\end{proposition}

In the language of Section~\ref{sec:glance}, $(P,D,(u,v))$ is a gap datum of sizes $(d+1,d)$ with trace exceeding $d$. It is used in the proof of Theorem~\ref{thm:main}.

\begin{proof}
\emph{Plan.} Pick one nonzero coefficient of the defect $C=I_d-YX$. Build a rational projection $p$ that annihilates, by sandwiching, every other group element in that entry of $C$. Then $P=\diag(I_d,p)$, and the factors $u,v$ are $X,Y$ extended by one row and one column manufactured from $C$ and $\bar p$. As in the proof of Lemma~\ref{lem:binary}, $C$ satisfies Equation~\eqref{eq:defect}, which only needs $XY=I_d$.

\emph{Step 1 (choose a coefficient).} Since $C\ne0$, some entry $C_{ij}$ has a group element $h_0$ with nonzero coefficient. That coefficient is $1$, because the field is $\F_2$. Multiplying on the left by $h_0^{-1}$ gives
\begin{equation}\label{eq:support}
 h_0^{-1}C_{ij}=1+\sum_{s\in S}s,
 \qquad S:=\{h_0^{-1}h: h\ne h_0\text{ in the support of }C_{ij}\}.
\end{equation}
$S$ is a finite subset of $H_2\setminus\{1\}$, possibly empty.

\emph{Step 2 (the projection $p$).} In $\Q[G_D]$ put
\[
 p_h=\frac{1+z_h+z_h^2}{3}\quad(h\in H_2),\qquad
 p=p_1\prod_{s\in S}(1-p_s).
\]
\begin{itemize}[leftmargin=2em]
\item $p_h^2=p_h$: the nine products in $(1+z_h+z_h^2)^2$ give each element of $C_3$ three times, so $(1+z_h+z_h^2)^2=3(1+z_h+z_h^2)$.
\item $p_h^*=p_h$: inversion permutes $\{1,z_h,z_h^2\}$.
\item All $p_h$ commute, because all $z_h$ commute.
\end{itemize}
A product of commuting self-adjoint idempotents is a self-adjoint idempotent, so $p=p^*=p^2$. Each of the $1+|S|$ factors of $p$ has denominator $3$, so $3^{1+|S|}p\in\Z[G_D]$.

\emph{Step 3 ($psp=0$ for $s\in S$).} Conjugation by $s$ is a ring automorphism of $\Q[G_D]$ with $sz_hs^{-1}=z_{sh}$, hence $sp_hs^{-1}=p_{sh}$ and
\[
 sps^{-1}=p_{s}\prod_{t\in S}(1-p_{st}).
\]
So $sps^{-1}$ has the factor $p_s$, while $p$ has the factor $1-p_s$. All factors commute and $(1-p_s)p_s=p_s-p_s^2=0$, so $p\,(sps^{-1})=0$. Therefore
\begin{equation}\label{eq:separation}
 psp=p\,(sps^{-1})\,s=0\qquad(s\in S).
\end{equation}
The argument does not need the label sets $\{1\}\cup S$ and $\{s\}\cup sS$ to be disjoint. Inverse pairs or finite-order elements in $S$ change nothing.

\emph{Step 4 (the trace of $p$).} The labels in $\{1\}\cup S$ are pairwise distinct because $1\notin S$. Expanding
\[
 p=3^{-(1+|S|)}\,(1+z_1+z_1^2)\prod_{s\in S}(2-z_s-z_s^2),
\]
a monomial $z_1^{a}\prod_sz_s^{b_s}$ with exponents in $\{0,1,2\}$ is the identity only if all exponents vanish, because distinct coordinates are independent. Hence
\begin{equation}\label{eq:theta}
 \tau_1(p)=[1]p=\frac13\prod_{s\in S}\frac23=\theta>0.
\end{equation}
For empty $S$ this reads $p=p_1$ and $\theta=1/3$.

\emph{Step 5 (reduction modulo two).} Since $3$ is odd, there is a ring homomorphism $\Z[1/3]\to\F_2$. Applied coefficientwise it gives a ring homomorphism $c\mapsto\bar c$ from $\Z[1/3][G_D]$ to $\F_2[G_D]$ fixing group elements. So $\bar p^2=\bar p$, and $\bar ps\bar p=0$ for $s\in S$ by Equation~\eqref{eq:separation}. With Equation~\eqref{eq:support},
\[
 \bar p\,h_0^{-1}C_{ij}\,\bar p
 =\bar p\Bigl(1+\sum_{s\in S}s\Bigr)\bar p=\bar p^2+\sum_{s\in S}\bar ps\bar p=\bar p.
\]

\emph{Step 6 (assembling $P,u,v$).} Let $f$ be the row of length $d$ with entry $\bar ph_0^{-1}$ at position $j$ and zero elsewhere. Let $w$ be the column of length $d$ with entry $\bar p$ at position $i$ and zero elsewhere. Then $fCw=\bar ph_0^{-1}C_{ij}\bar p=\bar p$ by Step~5. Define
\[
 P=\diag(I_d,p),\qquad
 u=\begin{pmatrix}X\\fC\end{pmatrix},\qquad
 v=\begin{pmatrix}Y&Cw\end{pmatrix}.
\]
The sizes are $(d+1)\times d$ and $d\times(d+1)$. By Equation~\eqref{eq:defect} ($XC=0$, $CY=0$, $C^2=C$) and $XY=I_d$,
\[
 uv=\begin{pmatrix}XY&XCw\\fCY&fC^2w\end{pmatrix}
   =\begin{pmatrix}I_d&0\\0&fCw\end{pmatrix}
   =\diag(I_d,\bar p)=\bar P.
\]
By Step~2, $P=P^*=P^2$ and $3^{1+|S|}P$ is integral. By Equation~\eqref{eq:theta}, $\tau_{d+1}(P)=d+\tau_1(p)=d+\theta$.
\end{proof}

The rational projection $P$ has trace greater than $d$, yet its binary reduction factors through $d$ columns. It is the rational projection, not its reduction, that carries the extra trace $\theta$.

\section{Lifting the factorization and building the integral matrix}
\label{sec:lifting}

From here until the end of Section~\ref{sec:analysis}, $G$ is an arbitrary group and $(P,D,(u,v))$ is a gap datum of sizes $(m,r)$ with $m,r\ge1$ (Section~\ref{sec:glance}).

\paragraph{Reductions.}
Let $q=2^k$, $k\ge1$, and $\mathcal R_q:=(\Z/q\Z)[G]$. Since $D$ is odd, it is a unit modulo $q$, so there is a ring homomorphism $\Z[1/D]\to\Z/q\Z$. For a matrix $\Phi$ over $\Z[1/D][G]$, write $\Phi_q$ for its coefficientwise image over $\mathcal R_q$ and $\bar\Phi$ for its image over $\F_2[G]$. On integral matrices these are the usual reductions, and reducing $\Phi_q$ further modulo two gives $\bar\Phi$.

\begin{lemma}\label{lem:lift}
Let $G$ be a group, $m,r\ge1$, $D$ an odd positive integer, and
$P\in\Mat_m(\Q[G])$ satisfy $P^2=P$ and
$DP\in\Mat_m(\Z[G])$. Suppose $\bar P=uv$ over $\F_2[G]$, where $u$ and $v$
have respective sizes $m\times r$ and $r\times m$. For $q=2^k$,
$k\ge1$, the reduction $e=P_q$ has a factorization
$e=u_qv_q$ of the same sizes over $(\Z/q\Z)[G]$.
\end{lemma}

Only idempotence of $P$ is used. The lemma feeds Lemma~\ref{lem:integral}.

\begin{proof}
\emph{Plan.} Lift $u,v$ arbitrarily. Their product, compressed by $e$, differs from $e$ by something even, and an element of the form $e-(\text{even})$ is invertible in the corner ring by a finite geometric series.

\emph{Step 1.} $e^2=e$, because reduction is a ring homomorphism. Choose arbitrary lifts $u_0,v_0$ of $u,v$ over $\mathcal R_q$ with finite supports. In the corner ring $e\Mat_m(\mathcal R_q)e$, whose identity is $e$, put
\[
 Z=eu_0v_0e,\qquad W=e-Z.
\]

\emph{Step 2 ($W^k=0$).} Modulo two, $Z$ becomes $\bar Puv\bar P=\bar P^3=\bar P$, which is also $e$ modulo two. So every coefficient of every entry of $W$ is even. Each coefficient of $W^k$ is a sum of products of $k$ such coefficients, hence divisible by $2^k=q$, so $W^k=0$ over $\mathcal R_q$. This is nilpotence of the even elements of the corner ring; the corner ring itself need not be nilpotent.

\emph{Step 3 (corner inverse).} Put $Q=e+W+\cdots+W^{k-1}$; for $k=1$ this is $e$. Since $We=eW=W$ and $Z=e-W$,
\[
 QZ=Q(e-W)=(e+W+\cdots+W^{k-1})-(W+\cdots+W^{k})=e-W^k=e.
\]
Thus $Q$ is an inverse of $Z$ in the corner ring, not in the whole matrix ring.

\emph{Step 4.} Set $u_q=Qeu_0$ and $v_q=v_0e$, of sizes $m\times r$ and $r\times m$. Then $u_qv_q=Q\,(eu_0v_0e)=QZ=e$. Only finite sums and products were used, so supports stay finite.
\end{proof}

\paragraph{Choice of $q$ and the target operator.}
Fix $q=2^k>1$ with $q\equiv1\pmod D$. Such $k\ge1$ exists because $2$ is a unit modulo the odd integer $D$: take $k$ a positive multiple of its multiplicative order. For $D=1$ any $k\ge1$ works. Let $u_q,v_q$ be as in Lemma~\ref{lem:lift}, and choose \emph{arbitrary} lifts $U\in\Mat_{m\times r}(\Z[G])$, $V\in\Mat_{r\times m}(\Z[G])$ of them. Define
\[
 J=-q^{-1}P-(I_m-P),\qquad B=\diag(qI_r,J).
\]
\emph{Motivation.} $B$ expands $r$ coordinate summands by $q$ and contracts the range of $P$ by $q^{-1}$, so its log-determinant should be $(r-\tau_m(P))\log q$. But $B$ is not integral. Prescribe upper-left block $qI_r$ and off-diagonal blocks $V,U$. Block elimination subtracts $U(qI_r)^{-1}V=UV/q$ from the lower-right block. For the Schur complement to be $J$, the lower-right block must be $UV/q+J$. The signs in $J$ are chosen so that this is $q^{-1}(UV-P)-(I_m-P)$, and $UV\equiv P\pmod q$ cancels the denominator.

\begin{lemma}\label{lem:integral}
Suppose $P=P^*=P^2\in\Mat_m(\Q[G])$, an odd $D>0$ clears its
denominators, and $\bar P=uv$ factors through $r\ge1$ columns.
Let $q=2^k>1$ satisfy $q\equiv1\pmod D$. Choose arbitrary
integral lifts $U,V$ of the factors supplied by Lemma~\ref{lem:lift}.
Then
\begin{equation}\label{eq:integral-matrix}
 A=\begin{pmatrix}
 qI_r&V\\
 U&q^{-1}\bigl(UV-P-q(I_m-P)\bigr)
 \end{pmatrix}
 \in\Mat_{r+m}(\Z[G])
\end{equation}
has an inverse in $\Mat_{r+m}(\Q[G])$. Moreover, for $x\in[0,1]$ and
\[
 L_x=\begin{pmatrix}I_r&0\\xU/q&I_m\end{pmatrix},\qquad
 R_x=\begin{pmatrix}I_r&xV/q\\0&I_m\end{pmatrix},
\]
one has
\begin{equation}\label{eq:path}
 M(x):=L_xBR_x=
 \begin{pmatrix}qI_r&xV\\xU&x^2UV/q+J\end{pmatrix},
 \qquad A=M(1).
\end{equation}
\end{lemma}

The path $M(x)$ from $B=M(0)$ to $A=M(1)$ is what Proposition~\ref{prop:criterion} uses.

\begin{proof}
\emph{Step 1 (integrality).} Write $P+q(I_m-P)=qI_m+(1-q)P$. It is integral because $D\mid(q-1)$ and $DP$ is integral. Its reduction modulo $q$ is $0+1\cdot P_q=e$. The reduction of $UV$ is $u_qv_q=e$ by Lemma~\ref{lem:lift}. Hence $UV-P-q(I_m-P)$ is an integral matrix all of whose coefficients are divisible by $q$, and the lower-right block of Equation~\eqref{eq:integral-matrix} is integral. The other blocks are integral by construction.

\emph{Step 2 ($B$ is invertible over $\Q[G]$).} Using $P^2=P$ and $P(I_m-P)=0$,
\[
 \bigl(q^{-1}P+(I_m-P)\bigr)\bigl(qP+(I_m-P)\bigr)=P+(I_m-P)=I_m,
\]
and likewise in the other order. So $J^{-1}=-qP-(I_m-P)$ and $B^{-1}=\diag(q^{-1}I_r,J^{-1})$.

\emph{Step 3 (the product).} Multiplying in order,
\[
 BR_x=\begin{pmatrix}qI_r&xV\\0&J\end{pmatrix},\qquad
 L_x(BR_x)=\begin{pmatrix}qI_r&xV\\xU&x^2UV/q+J\end{pmatrix}.
\]
At $x=1$ the lower-right block is $UV/q-q^{-1}P-(I_m-P)=q^{-1}\bigl(UV-P-q(I_m-P)\bigr)$, so $M(1)=A$.

\emph{Step 4 (inverse).} $L_x^{-1}$ and $R_x^{-1}$ are obtained by negating the off-diagonal block. Hence $A^{-1}=R_1^{-1}B^{-1}L_1^{-1}\in\Mat_{r+m}(\Q[G])$. More generally $M(x)^{-1}=R_x^{-1}B^{-1}L_x^{-1}$ has entries in $\R[G]$ and depends polynomially on $x$.
\end{proof}

No identities $PU=U$, $VP=V$ or $V=U^*$ are required of the lifts. The endpoint $A$ is integral; the intermediate $M(x)$ need not be.

\begin{example}[Sanity check with the trivial group]\label{ex:sanity}
Let $G=1$, $m=r=1$, $D=1$, $q=2$.
\begin{itemize}[leftmargin=2em]
\item $P=(1)$, $u=v=(1)$. Lemma~\ref{lem:lift} gives $u_q=v_q=1$ in $\Z/2\Z$. With lifts $U=3$, $V=1$: $A=\begin{pmatrix}2&1\\3&1\end{pmatrix}$, $|\det A|=1=q^{\,r-\tau_1(P)}=2^{0}$.
\item $P=(0)$, $u=v=(0)$. With lifts $U=2$, $V=4$: $A=\begin{pmatrix}2&4\\2&3\end{pmatrix}$, $|\det A|=2=2^{1-0}$.
\end{itemize}
Both agree with Proposition~\ref{prop:criterion} below. In both, $\tau_m(P)\le r$; Section~\ref{sec:projection} exists to produce a group where $\tau_m(P)>r$.
\end{example}

\section{The trace calculation}
\label{sec:analysis}

\paragraph{The algebra and its trace.}
The map $\alpha\mapsto T_\alpha$ on $\Mat_n(\C[G])$ is injective, since the entries of $\alpha$ can be read off from the vectors $T_\alpha\varepsilon_j$. We therefore write $\alpha$ for $T_\alpha$. Let $\mathcal A_n$ be the operator-norm closure of $\Mat_n(\C[G])$ in the bounded operators on $\ell^2(G)^n$; it is a unital $C^*$-algebra. The trace $\tau_n(T)=\sum_i\langle T\varepsilon_i,\varepsilon_i\rangle$ has these properties on $\mathcal A_n$.
\begin{enumerate}[label=(T\arabic*),leftmargin=3em]
\item $\tau_n$ is linear and $|\tau_n(T)|\le n\|T\|$, since each $|\langle T\varepsilon_i,\varepsilon_i\rangle|\le\|T\|$.
\item $\tau_n(T^*)=\overline{\tau_n(T)}$.
\item $\tau_n(TT')=\tau_n(T'T)$. For $\alpha,\beta\in\Mat_n(\C[G])$, Equation~\eqref{eq:trace-coefficient} gives
\[
 \tau_n(\alpha\beta)=\sum_{i,j}[1](\alpha_{ij}\beta_{ji})
 =\sum_{i,j}\sum_{g\in G}(\alpha_{ij})_g(\beta_{ji})_{g^{-1}},
\]
which is symmetric under $(\alpha,i,g)\leftrightarrow(\beta,j,g^{-1})$. The general case follows by density, (T1) and continuity of multiplication.
\end{enumerate}

\begin{lemma}\label{lem:logderivative}
For a continuously differentiable path
$M:[0,1]\to\mathcal A_n$ of invertible elements,
\[
 \frac{d}{dx}\tau_n\bigl(\log(M(x)^*M(x))\bigr)
 =2\operatorname{Re}\tau_n\bigl(M'(x)M(x)^{-1}\bigr).
\]
\end{lemma}

This is the tracial logarithmic derivative of~\cite[Lemma 2]{FK1952}, written with the unnormalized trace. It is used in Lemma~\ref{lem:shears}.

\begin{proof}
\emph{Plan.} Expand $\log N$, $N=M^*M$, in a norm-convergent power series, differentiate term by term, and use cyclicity (T3) to collapse the derivative of each power. No commutation of $N$ and $N'$ is assumed.

\emph{Step 1 (uniform bounds).} $N=M^*M$ is continuously differentiable with $N'=M'^*M+M^*M'$. Inversion is continuous on invertible elements, so $x\mapsto M(x)^{-1}$ is continuous on $[0,1]$. Put $c_1=1+\sup_x\|M(x)\|^2$ and $c_0=\inf_x\|M(x)^{-1}\|^{-2}>0$. For every vector $\xi$, $\langle N\xi,\xi\rangle=\|M\xi\|^2$ and $\|\xi\|\le\|M^{-1}\|\,\|M\xi\|$, so
\[
 c_0I\le N(x)\le c_1I\quad\text{for all }x,\qquad 0<c_0<c_1.
\]

\emph{Step 2 (series).} Put $\Lambda=I-N/c_1$. Then $0\le\Lambda\le(1-c_0/c_1)I$, so $\|\Lambda\|\le\eta:=1-c_0/c_1<1$. For $t\in[c_0,c_1]$ the scalar series $\log t=\log c_1-\sum_{j\ge1}(1-t/c_1)^j/j$ converges uniformly, so the continuous functional calculus gives
\[
 \log N=(\log c_1)I-\sum_{j\ge1}\frac{\Lambda^j}{j},
\]
norm-convergent uniformly in $x$.

\emph{Step 3 (termwise differentiation).} $\Lambda'=-N'/c_1$ is continuous, and $(\Lambda^j)'=\sum_{i=0}^{j-1}\Lambda^i\Lambda'\Lambda^{j-1-i}$, so
\[
 \left\|\frac{d}{dx}\frac{\Lambda^j}{j}\right\|
 \le\sup_{y\in[0,1]}\|\Lambda'(y)\|\,\eta^{j-1}.
\]
The series of derivatives therefore converges uniformly, and $\log N$ is continuously differentiable with derivative $-\sum_j(\Lambda^j)'/j$.

\emph{Step 4 (trace).} By (T3), $\tau_n(\Lambda^i\Lambda'\Lambda^{j-1-i})=\tau_n(\Lambda^{j-1}\Lambda')$ for each $i$, so $\tau_n((\Lambda^j)')/j=\tau_n(\Lambda^{j-1}\Lambda')$. Since $\tau_n$ is continuous and linear by (T1), and $\sum_{j\ge1}\Lambda^{j-1}=(I-\Lambda)^{-1}=c_1N^{-1}$,
\[
 \frac{d}{dx}\tau_n(\log N)
 =-\sum_{j\ge1}\tau_n(\Lambda^{j-1}\Lambda')
 =\frac1{c_1}\sum_{j\ge1}\tau_n(\Lambda^{j-1}N')
 =\tau_n(N^{-1}N').
\]

\emph{Step 5 (real part).} Insert $N'=M'^*M+M^*M'$ and $N^{-1}=M^{-1}(M^*)^{-1}$. By (T3),
\[
 \tau_n(N^{-1}M^*M')=\tau_n(M^{-1}M')=\tau_n(M'M^{-1}).
\]
By (T3), $MN^{-1}=(M^*)^{-1}$ and (T2),
\[
 \tau_n(N^{-1}M'^*M)=\tau_n(MN^{-1}M'^*)=\tau_n\bigl((M'M^{-1})^*\bigr)
 =\overline{\tau_n(M'M^{-1})}.
\]
The sum is $2\operatorname{Re}\tau_n(M'M^{-1})$.
\end{proof}

\begin{lemma}\label{lem:shears}
Let $r,m\ge1$ and let $B\in\mathcal A_{r+m}$ be invertible.
Let $\Gamma,\Delta$ have respective sizes $m\times r$ and $r\times m$ over
the operator-norm closure of $\C[G]$. For $x\in[0,1]$, let
\[
 L_x=\begin{pmatrix}I_r&0\\x\Gamma&I_m\end{pmatrix},\qquad
 R_x=\begin{pmatrix}I_r&x\Delta\\0&I_m\end{pmatrix}.
\]
Then $\tau_{r+m}(\log((L_xBR_x)^*(L_xBR_x)))$ is constant in $x$.
\end{lemma}

Here $B$ is any invertible element. Proposition~\ref{prop:criterion} applies the lemma with the $B$ of Section~\ref{sec:lifting}, $\Gamma=U/q$ and $\Delta=V/q$.

\begin{proof}
Write $M=L_xBR_x$. It is polynomial in $x$ and invertible, with $M^{-1}=R_x^{-1}B^{-1}L_x^{-1}$, where the shear inverses negate the off-diagonal blocks. So Lemma~\ref{lem:logderivative} applies. By the product rule, $M'=L_x'BR_x+L_xBR_x'$, hence
\[
 M'M^{-1}=L_x'L_x^{-1}
 +(L_xB)(R_x'R_x^{-1})(L_xB)^{-1}.
\]
Now
\[
 L_x'L_x^{-1}=\begin{pmatrix}0&0\\\Gamma&0\end{pmatrix}
 \begin{pmatrix}I_r&0\\-x\Gamma&I_m\end{pmatrix}
 =\begin{pmatrix}0&0\\\Gamma&0\end{pmatrix},\qquad
 R_x'R_x^{-1}=\begin{pmatrix}0&\Delta\\0&0\end{pmatrix}.
\]
Both have zero diagonal blocks, so trace zero. By (T3) the conjugate $(L_xB)(R_x'R_x^{-1})(L_xB)^{-1}$ also has trace zero. Thus $\tau_{r+m}(M'M^{-1})=0$, and Lemma~\ref{lem:logderivative} gives derivative zero.
\end{proof}

\begin{proposition}\label{prop:criterion}
Let $G$ be a discrete group, $m,r\ge1$, and let
$P=P^*=P^2\in\Mat_m(\Q[G])$. Suppose an odd $D>0$ clears its
denominators and $\bar P=uv$ with
$u\in\Mat_{m,r}(\F_2[G])$ and
$v\in\Mat_{r,m}(\F_2[G])$. For every $q=2^k>1$ congruent to one
modulo $D$, the matrix $A$ of Lemma~\ref{lem:integral} satisfies
\begin{align}
 \tau_{r+m}(\log(T_A^*T_A))
   &=2(r-\tau_m(P))\log q,\label{eq:criterion-trace}\\
 \log\det_{\N(G)}(T_A)&=(r-\tau_m(P))\log q.
 \label{eq:criterion-det}
\end{align}
It is integral and rationally invertible, and its positive-square
spectral measure has support in a compact subinterval of $(0,\infty)$. Consequently
\begin{equation}\label{eq:criterion-integral}
 \int_{(0,\infty)}\log t\,d\mu_A(t)
 =\frac{2(r-\tau_m(P))\log q}{r+m}.
\end{equation}
\end{proposition}

\begin{proof}
\emph{Step 1.} Integrality and the inverse over $\Q[G]$ are Lemma~\ref{lem:integral}.

\emph{Step 2 (reduce to $B$).} By Equation~\eqref{eq:path}, $A=M(1)$ and $B=M(0)$ for the path $M(x)=L_xBR_x$. $B$ is invertible in $\mathcal A_{r+m}$ by Step~2 of Lemma~\ref{lem:integral}. Lemma~\ref{lem:shears} with $\Gamma=U/q$, $\Delta=V/q$ gives
\[
 \tau_{r+m}(\log(A^*A))=\tau_{r+m}(\log(B^*B)).
\]

\emph{Step 3 (compute at $B$).} Since $P^*=P$ and $q$ is real, $J^*=J$. Since $P(I_m-P)=0$, $J^2=q^{-2}P+(I_m-P)$. So
\[
 B^*B=\diag\bigl(q^2I_r,\ q^{-2}P+I_m-P\bigr).
\]
This operator is $q^2$ on the first $r$ summands, $q^{-2}$ on the range of the orthogonal projection $P$ in the last $m$ summands, and $1$ on the orthogonal complement of that range. These three subspaces are mutually orthogonal and span, so $\log(B^*B)$ acts on them by $2\log q$, $-2\log q$ and $0$:
\[
 \log(B^*B)=2\log q\,\diag(I_r,-P),\qquad
 \tau_{r+m}(\log(B^*B))=2(r-\tau_m(P))\log q.
\]
With Step~2 this is Equation~\eqref{eq:criterion-trace}. Half of it is Equation~\eqref{eq:criterion-det}, by Equation~\eqref{eq:normalization}.

\emph{Step 4 (spectral support).} $T_A^{-1}=T_{A^{-1}}$ is bounded, since $A^{-1}$ is a finite group-ring matrix. So
\[
 \|T_A^{-1}\|^{-2}I\le T_A^*T_A\le\|T_A\|^2I,
\]
and $\mu_A$ is supported in $[\|T_A^{-1}\|^{-2},\|T_A\|^2]\subset(0,\infty)$. Equation~\eqref{eq:normalization} with $n=r+m$ turns Equation~\eqref{eq:criterion-trace} into Equation~\eqref{eq:criterion-integral}.
\end{proof}

\section{Assembly and exact determinant}
\label{sec:assembly}

\begin{proof}[Proof of Theorem~\ref{thm:main}]
\emph{Step 1.} Take $K_2,H_2,a_2,b_2$ from Theorem~\ref{thm:input} and put $d=[K_2:\F_2]\ge1$. Lemma~\ref{lem:binary} gives $X,Y\in\Mat_d(\F_2[H_2])$ with $XY=I_d\ne YX$.

\emph{Step 2.} Proposition~\ref{prop:separation} gives the group $G_D=(\bigoplus_{H_2}C_3)\rtimes H_2$, a finite set $S$, and an orthogonal projection $P\in\Mat_{d+1}(\Z[1/3][G_D])\subseteq\Mat_{d+1}(\Q[G_D])$ with
\[
 \tau_{d+1}(P)=d+\theta,\qquad\theta=\frac13\left(\frac23\right)^{|S|}>0,
 \qquad D=3^{1+|S|},\qquad\bar P=uv,
\]
where $u,v$ have sizes $(d+1)\times d$ and $d\times(d+1)$, and the odd integer $D$ clears the coefficients of $P$.

\emph{Step 3.} Choose $k\ge1$ with $q=2^k\equiv1\pmod D$. Proposition~\ref{prop:criterion}, with $r=d$ and $m=d+1$, produces $A\in\Mat_{2d+1}(\Z[G_D])$, invertible over $\Q[G_D]$, with
\[
 \log\det_{\N(G_D)}(T_A)=(d-(d+\theta))\log q=-\theta\log q.
\]
Since $q>1$ and $\theta>0$,
\begin{equation}\label{eq:final-det}
 0<\det_{\N(G_D)}(T_A)=q^{-\theta}<1.
\end{equation}

\emph{Step 4.} With $n=2d+1$, Equation~\eqref{eq:criterion-integral} reads
\begin{equation}\label{eq:final-integral}
 \int_{(0,\infty)}\log t\,d\mu_A(t)
 =-\frac{2\theta\log q}{2d+1}<0,
\end{equation}
which is finite by the spectral bounds in Proposition~\ref{prop:criterion}.

\emph{Step 5.} A finite generating set of $H_2$ together with $z_1$ generates $G_D$, since $z_h=hz_1h^{-1}$ for every $h\in H_2$. Thus $G_D$ is finitely generated.
\end{proof}

The target group $G_D$ differs from the companion's group $H_2$ and contains the cyclic coordinates of order three. The example has a finite negative logarithmic determinant and no zero spectral mass. It is a failure of the asserted lower bound, not of determinant class.

\section{Context}
\label{sec:context}

Nothing in this section is used in the proof. The statements about the literature are reported from the source.

\begin{itemize}[leftmargin=2em]
\item Determinant bounds are part of the theory of $L^2$-invariants \cite{Luck2002}. Approximation and closure arguments established such bounds for large classes of groups, including amenable groups \cite[Theorem 1.21, with its erratum]{Schick2001}. Elek and Szab\'o proved the determinant bound for sofic groups \cite[Theorem 5]{ElekSzabo2005}, using finite graph approximations and the integrality of products of nonzero eigenvalues of positive integral matrices. Applying that theorem to the positive integral matrix $A^*A$ shows that $G_D$ is not sofic. This is an external consequence, not an ingredient in constructing $A$.
\item The lower-bound question is distinct from determinant \emph{approximation}. Kammeyer's counterexample to a finite-quotient limiting formula has determinant $2$ over the discrete Heisenberg group, while its principal congruence quotients of levels $3^i$, $i\ge1$, have determinants at most $5^{1/3}$ \cite[Theorem 2 and Section 2]{Kammeyer2026}. That discrepancy does not violate the lower bound one.
\item The projections $p_h$ of Proposition~\ref{prop:separation} are related to the commuting wreath-product projections of Dicks and Schick \cite[Section 3]{DicksSchick2002} (compare \cite[Lemmas 3.1 and 3.4]{DicksSchick2002}) and of Lehner, Neuhauser and Woess~\cite[Section 2]{LehnerNeuhauserWoess2008}. Here they separate one coefficient of a binary defect; no spectral-distribution theorem is imported from those works.
\item Lemma~\ref{lem:logderivative} is the classical tracial logarithmic derivative \cite[Lemma 2]{FK1952}. It is also present in the path-integral determinant viewpoint of de la Harpe and Skandalis \cite[Lemma 1 and Proposition 2(d)]{deLaHarpeSkandalis1984}. Its proof and its application to block shears are given in full above.
\end{itemize}

\section{Verification checklist}
\label{sec:checklist}

One step is marked unverified: the external input, Theorem~\ref{thm:input}. The load-bearing points, in the order a sceptical reader should check them, are these.

\begin{enumerate}[leftmargin=2em]
\item \textbf{The external input.} Theorem~\ref{thm:input} is quoted from \cite{CharacteristicTwo} and is not checked here; see Remark~\ref{rem:input}. Everything else is proved in this document.
\item \textbf{Coefficient separation.} $psp=0$ for $s\in S$, and $[1]p=\theta$ (Proposition~\ref{prop:separation}, Steps~3 and~4). Check that the conjugation rule $sp_hs^{-1}=p_{sh}$ matches the left-translation action, and that the labels $\{1\}\cup S$ are distinct.
\item \textbf{The binary factorization.} $uv=\bar P$ (Proposition~\ref{prop:separation}, Steps~5 and~6). It uses $XC=0$, $CY=0$, $C^2=C$ from Lemma~\ref{lem:binary} and $\bar ph_0^{-1}C_{ij}\bar p=\bar p$. The coefficient at $h_0$ being exactly $1$ is where characteristic two with field $\F_2$ is used.
\item \textbf{The lift modulo $q$.} $W$ is even, $W^k=0$ and $QZ=e$ in the corner ring (Lemma~\ref{lem:lift}, Steps~2 and~3).
\item \textbf{Integrality of $A$.} The lower-right block needs both $q\equiv1\pmod D$ and $UV\equiv P\pmod q$ (Lemma~\ref{lem:integral}, Step~1).
\item \textbf{The shear factorization.} $L_xBR_x$ equals the displayed matrix and $A=M(1)$ (Lemma~\ref{lem:integral}, Step~3). Example~\ref{ex:sanity} tests this numerically.
\item \textbf{The logarithmic derivative.} Termwise differentiation and the use of cyclicity in place of commutativity (Lemma~\ref{lem:logderivative}, Steps~3 to~5); the trace-zero computation (Lemma~\ref{lem:shears}).
\item \textbf{The value at $B$ and the normalizations.} $\log(B^*B)=2\log q\,\diag(I_r,-P)$ (Proposition~\ref{prop:criterion}, Step~3); the factors $\tfrac12$ and $n=r+m$ in Equations~\eqref{eq:normalization} and~\eqref{eq:criterion-integral}; the final substitution $r=d$, $\tau_{d+1}(P)=d+\theta$ (proof of Theorem~\ref{thm:main}, Steps~3 and~4).
\end{enumerate}

\begingroup
\raggedright
\bibliographystyle{plain}
\bibliography{references}
\endgroup
\end{document}
